% !TeX program = xelatex
% UTF-8. Complete transcription of Notebook/vascular diseases of brain.
% Source photographs, in chronological filename order, are preserved unchanged.
% pic01.jpg = 20261012_210829985.jpg
% pic02.jpg = 20261012_210835939.jpg
% pic03.jpg = 20261012_210840563.jpg
% pic04.jpg = 20261012_210846492.jpg
% pic05.jpg = 20261012_210852254.jpg
% pic06.jpg = 20261012_210857772.jpg
% pic07.jpg = 20261012_210902936.jpg
% pic08.jpg = 20261012_210909183.jpg
% pic09.jpg = 20261012_210914594.jpg
% pic10.jpg = 20261012_210919876.jpg
% pic11.jpg = 20261012_210927123.jpg
% pic12.jpg = 20261012_210932548.jpg
% Medical statements and abbreviations are transcribed as written, not corrected.
\ifdefined\settextfont
  \def\VascularBrainStandalone{0}
\else
  \def\VascularBrainStandalone{1}
  \documentclass[11pt,a4paper]{article}
  \usepackage[margin=2cm]{geometry}
  \usepackage{amsmath,amssymb,enumitem}
  \usepackage{xepersian}
  \IfFontExistsTF{Amiri}{\settextfont{Amiri}}{\settextfont{FreeSerif}}
  \setlatintextfont{FreeSerif}
  \setlength{\emergencystretch}{2em}
  \setlist[itemize]{itemsep=3pt,topsep=4pt,parsep=0pt}
\fi
\providecommand{\en}[1]{\lr{#1}}
\newcommand{\vbarrow}{\ensuremath{\leftarrow}}
\newcommand{\vbpage}[1]{\clearpage\section*{بیماری‌های عروق مغزی — تصویر #1}
  {\footnotesize\noindent\en{Subject:} ۴ ـ بیماری‌های عروق مغزی\par
  \noindent\en{Date: \quad / \quad /}\par
  \noindent\en{Sat.\quad Sun.\quad Mon.\quad Tue.\quad Thu.\quad Wed.\quad Fri.}\par}
  \medskip}
\ifnum\VascularBrainStandalone=1 \begin{document}\fi

\vbpage{۱}
\subsection*{سکتهٔ مغزی}
\begin{itemize}
\item سکتهٔ مغزی \vbarrow{} آتروترومبوتیک، آمبولیک، لاکونار.
\item ویژگی‌های اصلی سکتهٔ مغزی:
  \begin{itemize}
  \item علائم نورولوژیک فقدان دارند [خوانش نامطمئنِ واژهٔ «فقدان»].
  \item علائم نورولوژیک فوکال است.
  \item علائم نورولوژیک به صورت حاد ایجاد شده و طی زمان سیر بدترشونده دارند.
  \end{itemize}
\item \en{RF} سکتهٔ مغزی \vbarrow{} \en{HTN, HLD, DM}، چاقی، سیگار، مصرف داروها، آبستنی و بیماری دریچه.
\item سن (مهم‌ترین)، \en{F < M}، نژاد سیاه‌پوستان.
\item بیماری‌های ارثی: کمبود فاکتورهای ضدانعقاد و یا افزایش هموسیستئین.
\end{itemize}

\subsection*{سکتهٔ مغزی آتروترومبوتیک}
\begin{itemize}
\item ایجاد علائم نسبت به فرم آمبولیک تدریجی‌تر است.
\item علائم در حین خواب رخ می‌دهند.
\item معمولاً سیر کند دارد و \en{TIA}هایی رخ می‌دهند که علائم مشابهی دارند.
\end{itemize}

\subsection*{سکتهٔ مغزی آمبولیک}
\begin{itemize}
\item منشأ: قلب، قوس آئورت، دیسکسیون عروق گردن.
\item در جوان‌ها: ترومای گردن.
\item علائم به صورت ناگهانی و کامل شروع می‌شوند.
\item اکثراً عروق اصلی درگیر هستند.
\item همی‌پلژی، آفازی، اختلال بینایی به صورت همونویمس آنون [خوانش نامطمئنِ اصطلاح دست‌نویس].
\end{itemize}

\subsection*{سایر علل سکتهٔ مغزی ایسکمیک}
واسکولیت‌ها، پلی‌سایتمی، ترومبوسیتوز، میگرن، آمبولی پارادوکس.
\par\hfill{\footnotesize شمارهٔ دست‌نویس صفحه: ۱}

\vbpage{۲}
\subsection*{سکتهٔ مغزی لاکونار}
\begin{itemize}
\item انسداد عروق ریز \en{$< 200\,\mathrm{mm}$} (واحد در دست‌نویس: \en{mm})؛ به دلیل لیپوهیالینوز \vbarrow{} هیپرتنشن.
\item \en{RF: HTN, DM}.
\item همی‌پارزی خالص (بدون علائم حسی)، اختلال همی‌سنسیوری خالص.
\item آتاکسی و همی‌پارزی، دیس‌آرتری و کلَمزی دست.
\end{itemize}
\subsection*{یافته‌های تصویربرداری در لاکونار}
\begin{itemize}
\item آنژیوگرافی معمولاً \en{NL}، چون عروق کوچک درگیر هستند.
\item \en{CSF} آنالیز \en{NL}.
\item \en{CT} و \en{MRI}، \en{NL}.
\end{itemize}

\subsection*{تغذیهٔ رگ درگیر؛ گردش خون قدامی مغز}
\begin{itemize}
\item کاروتید داخلی: \en{$2/3$} قدامی نیمکره‌ها و بخش‌هایی از دیانسفال.
\item شاخه‌های کاروتید داخلی:
  \begin{itemize}
  \item \en{ACA}: سطح داخلی نیمکره‌های مغزی.
  \item \en{MCA}: ادامهٔ کاروتید داخلی است / سطح خارجی نیمکره‌های مغزی.
  \item افتالمیک: خون‌رسانی شبکیه.
  \end{itemize}
\end{itemize}
\subsection*{علائم سکتهٔ ایسکمیک در مسیر شریان کاروتید داخلی}
\begin{itemize}
\item آموروزیس فوگاکس: کوری موقت یک چشم به صورت \en{TIA}.
\item همی‌پارزی؛ همی‌آنوپی.
\item اختلال همی‌سنسیوری.
\item آفازی و دیس‌فازی.
\item \en{Neglect}.
\end{itemize}
\par\hfill{\footnotesize شمارهٔ دست‌نویس صفحه: ۲}

\vbpage{۳}
\subsection*{علائم انفارکتوس شریان مغزی قدامی (\en{ACA})}
\begin{itemize}
\item ضعف حرکتی (پارزی) + بی‌اختیاری ادراری.
\item پارزی بیشتر در اندام تحتانی است / در اغلب موارد صورت درگیر نیست.
\item بی‌اختیاری ادرار و مدفوع معمولاً موقتی می‌باشد.
\end{itemize}

\subsection*{انفارکتوس \en{MCA}}
\begin{itemize}
\item ضعف حرکتی + آفازی + همی‌آنوپی.
\item پارزی معمولاً در اندام فوقانی و صورت.
\item آفازی در درگیری نیمکرهٔ غالب مشاهده می‌شود:
  \begin{itemize}
  \item درگیری تنهٔ \en{MCA} \vbarrow{} آفازی گلوبال.
  \item درگیری شاخهٔ فوقانی \en{MCA} \vbarrow{} آفازی بروکا.
  \item درگیری شاخهٔ تحتانی \en{MCA} \vbarrow{} آفازی ورنیکه.
  \end{itemize}
\item همی‌نوپسی هم‌آنوپی به دلیل درگیری \en{Optic Radiation} رخ می‌دهد.
\item در انفارکتوس \en{ACA} و \en{MCA} اختلالات حس سطحی (درد، حرارت، لمس، فشار) با توزیع مشابه اختلالات حرکتی مشاهده می‌گردد.
\item در انفارکتوس \en{MCA} علاوه بر حس سطحی، درگیری حس کورتیکال و حس عمقی نیز رخ می‌دهد؛ به صورت \en{Pseudoathetosis} است.
\end{itemize}

\subsection*{گردش خون خلفی}
\begin{itemize}
\item ۲ شریان ورتبرال؛ در اثر اتصال آن‌ها تنهٔ بازیلار تشکیل می‌شود.
\item ساقهٔ مغز، قسمت‌های تحتانی دیانسفال، \en{$1/3$} خلفی نیمکره‌ها.
\item شاخه‌های شریان بازیلار:
  \begin{itemize}
  \item \en{PCA}: سطوح خارجی و داخلی لوب اکسیپیتال و سطح داخلی لوب تمپورال.
  \end{itemize}
\end{itemize}
\par\hfill{\footnotesize شمارهٔ دست‌نویس صفحه: ۳}

\vbpage{۴}
\subsection*{شاخه‌های شریان بازیلار (ادامه)}
\begin{itemize}
\item شریان مخچه‌ای تحتانی خلفی.
\item شریان مخچه‌ای تحتانی قدامی.
\item شریان مخچه‌ای فوقانی.
\end{itemize}
\subsection*{علائم انفارکتوس شریان ورتبروبازیلار}
به علت درگیری ساقهٔ مغز (\en{brainstem}):
\begin{itemize}
\item کاهش سطح هوشیاری و کما.
\item افتالموپلژی.
\item ضعف حرکتی (پارزی): کوادری‌پارزی، همی‌پارزی.
\item فلج متقاطع: اعصاب کرانیال یک سمت، اندام‌های سمت مقابل.
\item سرگیجه، تهوع و استفراغ: علائم درگیری مخچه.
\item دیس‌آرتری و دیس‌فونی.
\item اختلال تعادل؛ آتاکسی.
\item کوری کورتیکال؛ همی‌آنوپی.
\item درگیری اعصاب کرانیال: دوبینی، فلج فاسیال، دیس‌فاژی، سرگیجه، اختلال شنوایی.
\end{itemize}
\subsection*{علائم انفارکتوس \en{PCA}}
\begin{itemize}
\item همی‌نوپسی هم‌آنوپی (شایع‌ترین).
\item ممکن است دید ماکولار حفظ شود: \en{PCA + MCA}.
\item اختلال حسی، اختلال در خواندن (\en{Alexia})، کنفیوژن.
\item اختلال حافظه.
\item آفازی آنومیک.
\end{itemize}
\subsection*{سکتهٔ هموراژیک}
خونریزی داخل بافت مغز و خارج بافت مغز.
\par\hfill{\footnotesize شمارهٔ دست‌نویس صفحه: ۴}

\vbpage{۵}
\subsection*{خونریزی داخل بافت مغز (\en{ICH})}
\begin{itemize}
\item پارگی عروق ریز شایع‌ترین علت.
\item همان محل شایع مشابه انفارکتوس لاکونار است.
\item \en{RF: HTN}.
\item علائم بالینی: سردرد، استفراغ، کاهش سطح هوشیاری، تشنج.
\end{itemize}
\subsection*{خونریزی خارج بافت مغز}
ساب‌آراکنوئید (\en{SAH})، ساب‌دورال، اپی‌دورال.
\subsubsection*{\en{SAH}}
\begin{itemize}
\item پارگی آنوریسم‌های \en{Berry} در مغز (شایع).
\item تروما (شایع).
\item کوآگولوپاتی (ناشایع).
\item پارگی آنوریسم مایکوتیک (ناشایع).
\item علائم: سردرد ناگهانی شدید، ریجید گردن، کنفیوژن، خونریزی در هیالوئید ـ شبکیه.
\end{itemize}
\subsubsection*{\en{SDH}}
\begin{itemize}
\item تروما و اختلالات انعقادی.
\item در تصویربرداری هلالی با تحدب به داخل مغز.
\end{itemize}
\subsubsection*{\en{EDH}}
\begin{itemize}
\item تروما شایع‌ترین علت آن است.
\item در تصویربرداری عدسی‌شکل با تحدب به سمت داخل مغز.
\end{itemize}
\subsection*{آزمایش‌های مربوط به انفارکتوس مغزی؛ مظنون به آن}
\begin{latin}
\noindent CBC-diff, ESR, BS, Total Chol, LDL, Electrolytes, PT, PTT\par
\noindent Kidney and Liver function tests
\end{latin}
\begin{itemize}
\item تصویربرداری \vbarrow{} \en{CT} بدون کنتراست.
\end{itemize}
\par\hfill{\footnotesize شمارهٔ دست‌نویس صفحه: ۵}

\vbpage{۶}
\subsection*{دلایل انجام \en{CT}}
\begin{itemize}
\item افتراق بیماری‌های مشابه سکتهٔ مغزی.
\item افتراق خونریزی از انفارکتوس:
  \begin{itemize}
  \item خونریزی: هایپردنس.
  \item انفارکت: هایپودنس در محدودهٔ رگ مسدود.
  \end{itemize}
\end{itemize}
\subsection*{ارزیابی قلبی بیمار جهت کشف منشأ آمبولی مغزی}
\begin{itemize}
\item \en{EKG} و آنزیم‌های قلبی.
\item به طور شایع \en{Stroke} در ۲ هفتهٔ اول بعد از \en{MI}، خصوصاً انفارکتوس قدامی قلب، رخ می‌دهد. در این‌ها \en{Stroke} اولین تظاهر \en{MI} است.
\item لوب اینسولا شایع‌ترین محل درگیر برای آمبولی‌های ناشی از آریتمی قلبی.
\item \en{TTE} در تمامی بیماران مبتلا به سکتهٔ مغزی الزامی است.
\item موارد مطرح‌کنندهٔ علل قلبی در \en{TTE}:
  \begin{itemize}
  \item \en{$EF < 15\%$}، بیماری‌های دریچه‌ای قلب، فیبری دهلیز چپ [خوانش نامطمئنِ واژهٔ «فیبری»].
  \item آکینزی دیوارهٔ بطن چپ، ترومبوز غیرکلسیفیهٔ بطن چپ.
  \end{itemize}
\item \en{TEE}.
\item در بیمار زیر ۴۵ سال که علت \en{Stroke} مشخص نشود، \en{TTE + TEE} انجام می‌شود.
\end{itemize}
\subsection*{سونوگرافی داپلر گردن}
\begin{itemize}
\item تشخیص تنگی و دایسکشن.
\item در صورت وجود یافتهٔ غیرطبیعی \vbarrow{} \en{CT} آنژیو، \en{MR} آنژیو، آنژیوگرافی.
\item در بیمار با سابقهٔ \en{TIA} باید جهت بررسی وجود تنگی عروق گردن \vbarrow{} سونو داپلر عروق گردن.
\end{itemize}
\subsection*{\en{TIA}}
\begin{itemize}
\item \en{low Risk}: احتمال \en{Stroke} طی ۳ ماه \en{$< 5\%$}.
\item \en{high Risk}: احتمال \en{Stroke} طی ۳ ماه \en{$\simeq 20\%$}.
\end{itemize}
\par\hfill{\footnotesize شمارهٔ دست‌نویس صفحه: ۶}

\vbpage{۷}
\subsection*{\en{TIA} و معیارهای \en{ABCD2}}
\begin{itemize}
\item اکثر \en{Stroke}ها ظرف ۴۸ ساعت اول پس از \en{TIA} رخ می‌دهند.
\item معیارهای \en{ABCD2}: حداکثر امتیاز ۷ / امتیاز \en{$\geq 5$} به معنای \en{high Risk} بوده؛ بستری و به عنوان \en{Stroke} در حال وقوع تلقی شود.
\item سن: بالای ۶۰: ۱ امتیاز / زیر ۶۰: ۰ امتیاز.
\item \en{BP}: \en{$SBP > 140$} یا \en{$DBP > 90$}: ۱ امتیاز.
\item مدت علائم: بیشتر از ۱ ساعت: ۲ امتیاز / ۱۰ دقیقه تا ۱ ساعت: ۱ امتیاز.
\item \en{DM}: دارد: ۱ امتیاز / ندارد: ۰ امتیاز.
\item تظاهر بالینی: همی‌پارزی: ۲ امتیاز / اختلال تکلم: ۱ امتیاز.
\end{itemize}

\subsection*{درمان دارویی در انفارکتوس مغزی}
\begin{itemize}
\item \en{rtPA}: ترومبولیتیک / درمان انتخابی انفارکتوس ایسکمیک طی ۴٫۵ ساعت اول.
\item میزان بهبودی سه‌ماهه را تا ۳۰\% افزایش می‌دهد.
\item \en{$0.9\,\mathrm{mg/kg}$}: ۱۰\% بولوس + ۹۰\% انفوزیون طی یک ساعت.
\item بهترین زمان طی ۳ ساعت اول شروع علائم است.
\item روش مصرف: \en{IV, IA}.
\item مانیتورینگ: بستری؛ طی ۲۴ ساعت اول مانیتور شود.
\item عارضهٔ اصلی: خونریزی؛ درمان: \en{FFP}، پلاکت، کرایو.
\end{itemize}

\subsection*{هپارین}
\begin{itemize}
\item در بیمار \en{Stroke} تجویز هپارین با دوز درمانی به صورت روتین توصیه نمی‌گردد.
\item اندیکاسیون هپارین با دوز درمانی:
  \begin{itemize}
  \item وجود بیماری قلبی \en{high Risk} برای وقوع آمبولی.
  \item بیماری‌های عروق بزرگ گردن مثل دایسکشن ورتبرال و کاروتید.
  \item شرایط هایپرکوآگولوپاتی.
  \item ترومبوز سینوس‌های وریدی مغز.
  \end{itemize}
\end{itemize}
\par\hfill{\footnotesize شمارهٔ دست‌نویس صفحه: ۷}

\vbpage{۸}
\subsection*{هپارین (ادامه)}
\begin{itemize}
\item هپارین همچنین برای پیشگیری از \en{DVT} نیز کاربرد دارد.
\item دوز درمانی هپارین: ۱۲۰۰ـ۱۵۰۰ واحد در ساعت.
\item چون نیمه‌عمر آن \en{2h} است، بهترین روش تجویز به کارگیری پمپ انفوزیون ثابت.
\item مهم‌ترین اندیکاسیون هپارین در درمان، فیبریلاسیون دهلیزی است.
\end{itemize}

\subsection*{داروهای ضدپلاکتی}
\begin{itemize}
\item داروی اصلی این دسته آسپرین است.
\item باید ظرف ۴۸ ساعت اول \en{Stroke} شروع شود.
\item در سکته‌های آتروترومبوتیک و لاکونار، از ایجاد ترومبوز مجدد جلوگیری می‌کند.
\item همچنین از بروز سکتهٔ قلبی و مغزی مجدد جلوگیری می‌کند.
\item در \en{Stroke} ایسکمیک اگر از ۴٫۵ ساعت گذشته باشد، باید آسپرین تجویز گردد.
\item بیماری در مسیر شریان یک طرف تنگی ناکامل و با علائم باشد \vbarrow{} فقط آسپرین.
\end{itemize}

\subsection*{استاتین‌ها}
\begin{itemize}
\item کاهش احتمال سکتهٔ مغزی.
\item کاهش کلسترول خون، اثر ضدالتهاب، پایدار نمودن پلاک آترومایی، تسریع در کاهش اندازهٔ پلاک آترومایی، بهبود پرفیوژن، پیشگیری از سکتهٔ قلبی و مغزی.
\end{itemize}
\subsection*{جراحی}
\begin{itemize}
\item در انفارکتوس‌های بدخیم ترانس‌تنتوریس [خوانش نامطمئنِ اصطلاح دست‌نویس] کمک‌کننده است.
\item انفارکتوس بدخیم: ادم شدید و هرنی‌شدن محتویات در پی انفارکتوس وسیع در مسیر شریان \en{MCA}.
\end{itemize}

\subsection*{درمان خونریزی‌های مغزی}
\begin{itemize}
\item اقدامات طبی و جراحی.
\item جراحی جهت تخلیهٔ هماتوم و پیشگیری از هرنی مغز.
\item قبل از جراحی باید علت خونریزی مشخص شود.
\item بیماری‌های زمینه‌ای مثل \en{HTN} و اختلالات انعقادی کنترل شوند.
\end{itemize}
\par\hfill{\footnotesize شمارهٔ دست‌نویس صفحه: ۸}

\vbpage{۹}
\subsection*{درمان خونریزی‌های مغزی (ادامه)}
\begin{itemize}
\item درمان طبی غالباً حمایتی است.
\item مانیتول الزامی است برای کاهش ادم مغزی و کاهش \en{ICP}.
\end{itemize}

\subsection*{ترومبوز سینوس‌ها و وریدهای مغزی (\en{CVST})}
\begin{itemize}
\item علل عفونی: استئومیلیت جمجمه، عفونت مجاور، در گوش داخلی، یا سینوس [خوانش نامطمئنِ عبارت پایانی].
\item علل غیرعفونی \vbarrow{} شایع‌ترین علل:
  \begin{itemize}
  \item هایپرکوآگولوپاتی، دهیدراتاسیون، مصرف \en{OCP}، حاملگی، دورهٔ نفاس، واسکولیت‌ها.
  \end{itemize}
\item در ۲۰\% موارد علت یافت نمی‌شود.
\item موارد غیرعفونی شایع‌ترین علت ترومبوز سینوس‌های ساژیتال با گسترش به سینوس‌های لترال و ژوگولار هستند.
\item علائم: سردرد شدید در رأس آنها (ممکن است تنها علامت باشد)، استفراغ، تاری دید گذرا، تشنج جنرالیزه یا فوکال، خواب‌آلودگی، کما.
\item ادم پاپی دوطرفه نیز شایع است.
\end{itemize}
\subsubsection*{ترومبوز سینوس کاورنوس}
\begin{itemize}
\item به دلیل مجاورت با اعصاب ۳، ۴ و ۶: پروپتوز، کموز، افتالموپلژی دردناک.
\end{itemize}
\subsubsection*{اگر ترومبوز سینوس کاورنوس خلفی ـ سینوس پتروزال تحتانی ـ گسترش یابد}
می‌تواند اعصاب ۹ و ۱۰ و ۱۱ را نیز درگیر کند.
\subsubsection*{ترومبوز سینوس پتروزال فوقانی}
\begin{itemize}
\item متعاقب \en{OM} یا ماستوئیدیت: فلج عصب ۵ و ۶.
\item در موارد کاورنوس درد، تب، سرگیجه، ادم پاپی دوطرفه (به دلیل \en{$\uparrow ICP$}).
\end{itemize}
\subsubsection*{تشخیص}
\begin{itemize}
\item بسیار دشوار است / اورژانس نورولوژی است / سریعاً کشنده.
\item علائم عفونی بالینی + تصویربرداری از مغز.
\item \en{CT} کمک‌کننده است ولی \en{MRV} و \en{MRI} حساس‌تر هستند.
\end{itemize}
\par\hfill{\footnotesize شمارهٔ دست‌نویس صفحه: ۹}

\vbpage{۱۰}
\subsection*{تصویربرداری ترومبوز وریدی مغز (ادامه)}
\begin{itemize}
\item ترومبوز و لخته‌ها در \en{$T_2$}: \en{Loss of the normal flow void}.
\item توقف جریان خون سینوس‌ها در \en{MRV} دیده می‌شود.
\item ترومبوز سینوس ساژیتال فوقانی در \en{CT} با کنتراست:
  \begin{itemize}
  \item عدم جذب مادهٔ حاجب در سینوس ساژیتال خلفی: \en{Empty Delta Sign}.
  \item انفارکتوس هموراژیک ناحیهٔ فرونتال یا پریتال پارامدین دوطرفه ممکن است دیده شود.
  \end{itemize}
\item ترومبوز سینوس لترال \vbarrow{} انفارکتوس هموراژیک، پارانشیم لوب تمپورال، ادم وازوژنیک.
\item بهترین روش تشخیص ترومبوز وریدهای عمقی مغز \vbarrow{} \en{MRI} با کنتراست.
\item تغییر سیگنال تالاموس دوطرفه به شکل دو لوبیا به هم چسبیده.
\end{itemize}

\subsection*{درمان}
\begin{itemize}
\item اصلاح عامل زمینه‌ای.
\item کاهش \en{ICP}؛ پروفیلاکسی تشنج.
\item تجویز هپارین \en{IV}، \en{LMWH SC}، سپس وارفارین.
\item تزریق موضعی ترومبولیتیک در سینوس وریدی درگیر: در صورت وجود کنترااندیکاسیون نسبت به تجویز هپارین، عدم پاسخ درمان، وخیم‌شدن بیماری.
\item گاهی جراحی \en{Decompressive}.
\end{itemize}
\subsection*{طول مدت درمان با آنتی‌کوآگولان}
\begin{itemize}
\item در موارد ایدیوپاتیک: ۳ـ۶ ماه.
\item در اختلالات هایپرکوآگولوپاتی هتروزیگوت: ۶ـ۱۲ ماه.
\item در اختلالات هایپرکوآگولوپاتی هموزیگوت: وارفارین تا پایان عمر.
\end{itemize}
\par\hfill{\footnotesize شمارهٔ دست‌نویس صفحه: ۱۰}

\vbpage{۱۱}
\subsection*{موارد درمان تا پایان عمر}
\begin{itemize}
\item \en{APS}.
\item تکرار و عود ترومبوز ورید مغزی.
\item وجود همزمان ترومبوآمبولی در دیگر نقاط بدن.
\end{itemize}

\subsection*{درمان ترومبوزهای عفونی}
\begin{itemize}
\item درمان اصلی تجویز زودهنگام \en{AB} است.
\item شایع‌ترین پاتوژن: استاف اورئوس.
\item درمان ایده‌آل:
  \begin{itemize}
  \item نفس‌سیلین \en{IV}، ۱٫۵ گرم هر ۴ ساعت.
  \item سفوتاکسیم \en{IV}، ۱٫۵ـ۲ گرم هر ۴ ساعت.
  \item مترونیدازول \en{IV}، \en{$15\,\mathrm{mg/kg}$} و سپس \en{$7.5\,\mathrm{mg/kg}$} هر ۶ ساعت.
  \end{itemize}
\item در صورت شک به \en{MRSA}: ونکومایسین جایگزین نفس‌سیلین.
\item درمان \en{AB} وریدی باید حداقل ۳ـ۴ هفته ادامه پیدا کند.
\end{itemize}

\subsection*{عوارض ترومبوز مغزی}
\begin{itemize}
\item انفارکتوس وریدی (وخیم‌ترین پیامد).
\item سودوتومور سربری.
\item سکته‌های هموراژیک غیرقابل توجیه.
\end{itemize}

\subsection*{انسفالوپاتی هایپرتنسیو (انسفالوپاتی خلفی برگشت‌پذیر)}
\begin{itemize}
\item افزایش سریع و شدید \en{BP}: \en{$\geq 180/120$}.
\item علل: \en{CKD}، تنگی شریان کلیوی، گلومرولونفریت حاد، مسمومیت حاد، فئوکروموسایتوم، سندرم کوشینگ، اکلامپسی.
\item داروهایی مثل آمینوفیلین یا فنیل‌افرین.
\end{itemize}
\par\hfill{\footnotesize شمارهٔ دست‌نویس صفحه: ۱۱}

\vbpage{۱۲}
\subsection*{انسفالوپاتی هایپرتنسیو (ادامه)}
\begin{itemize}
\item اکلامپسی نوعی انسفالوپاتی هایپرتنسیو است.
\item در زنان باردار، کودکان و بیماران با اختلالات کلیوی حاد، علائم انسفالوپاتی هایپرتنسیو در فشار خون‌های پایین‌تری آغاز می‌شود.
\item انسفالوپاتی هایپرتنسیو بیشتر در افراد جوان و میان‌سال مبتلا به \en{HTN}.
\item علائم به دنبال ۱۲ـ۴۸ ساعت بالا رفتن ناگهانی و مداوم فشار خون ایجاد می‌شود.
\item سردرد شدید (تظاهر اولیهٔ اکثر بیماران).
\item استفراغ جهنده.
\item گیجی شدید به علت \en{$\uparrow ICP$}.
\item تاری دید، نقص یک‌طرفهٔ میدان بینایی، کوررنگی، کوری کورتیکال.
\item در موارد پیشرفته: استوپور، کما، تشنج جنرالیزه.
\item همی‌پارزی و آفازی شایع نیستند.
\end{itemize}

\subsection*{یافته‌های تصویربرداری}
\begin{itemize}
\item افزایش دانسیته در \en{CT}، افزایش سیگنال \en{MRI} به دلیل تجمع مایع در خلف اکسیپوت.
\item در ابتدا یافته‌ها در \en{CT} و \en{MRI} با انفارکتوس مغزی و مناطق دمیلیناسیون اشتباه گرفته می‌شوند ولی به دنبال چند هفته یافته‌ها محو می‌گردند.
\end{itemize}
\begin{latin}
\noindent Posterior Reversible Encephalopathy Syndrome (PRES)
\end{latin}

\subsection*{درمان}
\begin{itemize}
\item کاهش صحیح و سریع فشار خون.
\item هدف \en{BP}: \en{$150/100\,\mathrm{mmHg}$}.
\item داروها:
  \begin{itemize}
  \item نیتروپروساید سدیم \en{IV}.
  \item \en{CCB} مثل نیفدیپین زیرزبانی.
  \item \en{BB} وریدی مثل لابتالول و اسمولول.
  \end{itemize}
\item داروی ضدتشنج در صورت تشنج پایدار.
\item اکلامپسی: سولفات منیزیم.
\item دگزا ۴ـ۶ \en{mg} هر ۶ ساعت اگر علائم \en{$\uparrow ICP$}.
\end{itemize}
\par\hfill{\footnotesize شمارهٔ دست‌نویس صفحه: ۱۲}

\ifnum\VascularBrainStandalone=1
  \def\vbfinish{\end{document}}
\else
  \let\vbfinish\relax
\fi
\vbfinish
